\subsection{Orthogonality, Radicals, and Nonsingular Metric Spaces}

Throughout this subsection, metric vector spaces are finite-dimensional over an arbitrary field $\mathbb{F}$.

\begin{definition} \label{def:metric-orthogonality-isotropy-radical} \hlt{Orthogonality, Isotropy, and Radical}\\
Let $\mathcal{V}$ be a metric vector space. Symmetric and skew-symmetric forms satisfy $(\forall\,\mathbf{x},\mathbf{y}\in\mathcal{V})\,[\innerprod{\mathbf{x},\mathbf{y}}=0\Longleftrightarrow\innerprod{\mathbf{y},\mathbf{x}}=0]$.\\
An alternating form is skew-symmetric by Proposition \ref{prop:form-type-relations-by-characteristic}(i); thus left and right orthogonality agree in every case.
\begin{enumerate}[label=(\roman*)]
\setlength{\itemsep}{0pt}
\item $\mathbf{x}\perp\mathbf{y}\Longleftrightarrow\innerprod{\mathbf{x},\mathbf{y}}=0$; then $\mathbf{x},\mathbf{y}$ are \hlt{orthogonal}. For $S\subseteq\mathcal{V}$, set $S^{\perp}\defeq\{\mathbf{v}\in\mathcal{V}:(\forall\,\mathbf{s}\in S)\,\innerprod{\mathbf{v},\mathbf{s}}=0\}$.
\item A subspace $\mathcal{S}\leq\mathcal{V}$ is \hlt{totally isotropic} or \hlt{totally degenerate} if $\mathcal{S}\subseteq\mathcal{S}^{\perp}$. It is \hlt{nonsingular} if $\mathcal{S}\cap\mathcal{S}^{\perp}=\{\mathbf{0}\}$.
\item A vector $\mathbf{x}\neq\mathbf{0}$ is \hlt{isotropic} or \hlt{null} if $\innerprod{\mathbf{x},\mathbf{x}}=0$, and \hlt{nonisotropic} otherwise.\\
The space $\mathcal{V}$ is \hlt{isotropic} if $(\exists\,\mathbf{x}\in\mathcal{V}\setminus\{\mathbf{0}\})\,\innerprod{\mathbf{x},\mathbf{x}}=0$. It is \hlt{anisotropic} if $(\forall\,\mathbf{x}\in\mathcal{V}\setminus\{\mathbf{0}\})\,\innerprod{\mathbf{x},\mathbf{x}}\neq0$.
\item A vector $\mathbf{x}$ is \hlt{degenerate} if $(\forall\,\mathbf{v}\in\mathcal{V})\,\innerprod{\mathbf{x},\mathbf{v}}=0$. The \hlt{radical} is $\operatorname{rad}(\mathcal{V})\defeq\mathcal{V}^{\perp}$.
\item The space $\mathcal{V}$ is \hlt{nonsingular} or \hlt{nondegenerate} if $\operatorname{rad}(\mathcal{V})=\{\mathbf{0}\}$, and \hlt{singular} otherwise.
\item The space $\mathcal{V}$ is \hlt{totally singular} or \hlt{totally degenerate} if $\operatorname{rad}(\mathcal{V})=\mathcal{V}$.
\end{enumerate}
\end{definition}

\begin{proposition} \label{prop:metric-radical-equals-bilinear-radicals} \hlt{Metric Radical and Perfection}\\
Let $\mathcal{V}$ be a metric vector space with form $B(\mathbf{x},\mathbf{y})=\innerprod{\mathbf{x},\mathbf{y}}$.
\begin{enumerate}[label=(\roman*)]
\setlength{\itemsep}{0pt}
\item $\operatorname{rad}_{L}(B)=\operatorname{rad}_{R}(B)=\operatorname{rad}(\mathcal{V})$.
\item $\mathcal{V}$ is nonsingular $\Longleftrightarrow B$ is perfect.
\end{enumerate}
\begin{proof}
\begin{enumerate}[label=(\roman*)]
\setlength{\itemsep}{0pt}
\item[$(i)$] The identity $\operatorname{rad}(\mathcal{V})=\operatorname{rad}_{L}(B)$ holds by the definitions of $\mathcal{V}^{\perp}$ and the left radical.\\
The equality $\operatorname{rad}_{L}(B)=\operatorname{rad}_{R}(B)$ holds by Proposition \ref{prop:form-type-relations-by-characteristic}(iv).
\item[$(ii)$] ($\Rightarrow$) Suppose $\mathcal{V}$ is nonsingular. Both form-induced maps are linear by Proposition \ref{prop:bilinear-forms-linearize-through-duals}(i).\\
Their kernels equal $\{\mathbf{0}\}$ by (i) and Proposition \ref{prop:bilinear-radicals-matrix-and-dual-maps}(i). Thus both maps are injective by Proposition \ref{prop:injective-iff-kernel-zero}.\\
Their domains and codomains have equal dimension by Theorem \ref{thm:dual-basis-and-dimension}(ii).\\
Hence both maps are surjective by Corollary \ref{cor:equal-dimension-kernel-range}(ii). They are isomorphisms by Proposition \ref{prop:linear-map-invertible-iff-bijective}, so $B$ is perfect.
\item[$(ii)$] ($\Leftarrow$) Suppose $B$ is perfect. Then $\Phi_B$ is an isomorphism by the definition of perfection. It is injective by Proposition \ref{prop:linear-map-invertible-iff-bijective}.\\
Hence $\ker(\Phi_B)=\{\mathbf{0}\}$ by Proposition \ref{prop:injective-iff-kernel-zero}. Thus $\operatorname{rad}(\mathcal{V})=\operatorname{rad}_{L}(B)=\ker(\Phi_B)=\{\mathbf{0}\}$ by (i) and Proposition \ref{prop:bilinear-radicals-matrix-and-dual-maps}(i).
\end{enumerate}
\hlt{Structural Commentary}: The common radical is the kernel obstruction; equal dimensions convert injectivity into perfection.\\
\hlt{Key Insight}: A metric form is perfect exactly when its radical vanishes.
\end{proof}
\end{proposition}

\begin{proposition} \label{prop:metric-radical-coordinate-criteria} \hlt{Coordinate Criteria for Nonsingularity}\\
Let $\mathcal{V}$ be a metric vector space with $n\defeq\dim(\mathcal{V})\geq1$. Let $\beta$ be an ordered basis, and set $M\defeq M_{\beta}$.
\begin{enumerate}[label=(\roman*)]
\setlength{\itemsep}{0pt}
\item $\{[\mathbf{x}]_{\beta}:\mathbf{x}\in\operatorname{rad}(\mathcal{V})\}=\ker(M^t)$.
\item $\dim\operatorname{rad}(\mathcal{V})=n-\operatorname{rank}(M)$.
\item $\mathcal{V}$ is nonsingular $\Longleftrightarrow M$ is invertible.
\end{enumerate}
\begin{proof}
\begin{enumerate}[label=(\roman*)]
\setlength{\itemsep}{0pt}
\item[$(i)$] For $\mathbf{x}\in\mathcal{V}$, set $X\defeq[\mathbf{x}]_{\beta}$. The relation $(\forall\,Y\in\mathbb{F}^n)(\exists\,\mathbf{y}\in\mathcal{V})\,Y=[\mathbf{y}]_{\beta}$ holds by Proposition \ref{prop:coordinate-map-isomorphism}(i).\\
Thus $\mathbf{x}\in\operatorname{rad}(\mathcal{V})\Longleftrightarrow(\forall\,Y\in\mathbb{F}^n)\,X^tMY=0$ by Proposition \ref{prop:metric-form-matrix-recovery-and-congruence}(i).\\
The equivalence $(\forall\,Y\in\mathbb{F}^n)\,X^tMY=0\Longleftrightarrow M^tX=0$ holds by evaluation at the standard coordinate columns. QED.
\item[$(ii)$] By (i) and Proposition \ref{prop:coordinate-map-isomorphism}(i), $\operatorname{rad}(\mathcal{V})\cong\ker(M^t)$. Hence their dimensions agree by Theorem \ref{thm:finite-dimensional-isomorphic-iff-same-dimension}.\\
Thus $\dim\operatorname{rad}(\mathcal{V})=n-\operatorname{rank}(M^t)$ by Corollary \ref{cor:matrix-rank-nullity}(i). Finally, $\operatorname{rank}(M^t)=\operatorname{rank}(M)$ by Theorem \ref{thm:row-rank-equals-column-rank}.
\item[$(iii)$] ($\Rightarrow$) Suppose $\mathcal{V}$ is nonsingular. Then $\operatorname{rank}(M)=n$ by (ii). Hence $M$ is invertible by Corollary \ref{cor:square-matrix-invertible-iff-full-rank}.
\item[$(iii)$] ($\Leftarrow$) Suppose $M$ is invertible. Then $\operatorname{rank}(M)=n$ by Corollary \ref{cor:square-matrix-invertible-iff-full-rank}.\\
Hence $\dim\operatorname{rad}(\mathcal{V})=0$ by (ii). Hence $\operatorname{rad}(\mathcal{V})=\{\mathbf{0}\}$ by Proposition \ref{prop:dimension-bounds-and-basis-criteria}(ii).
\end{enumerate}
\hlt{Structural Commentary}: Coordinate evaluation identifies the radical kernel; rank-nullity and full rank test nonsingularity.\\
\hlt{Key Insight}: The form is nonsingular exactly when its matrix has full rank.
\end{proof}
\end{proposition}

\begin{remark} \label{rem:subspace-radical-versus-ambient-complement} \hlt{Subspace Radicals}\\
For $\mathcal{S}\leq\mathcal{V}$, the restricted form has radical $\operatorname{rad}(\mathcal{S})=\mathcal{S}\cap\mathcal{S}^{\perp}$. In general, $\operatorname{rad}(\mathcal{S})\neq\mathcal{S}^{\perp}$ because $\mathcal{S}^{\perp}$ is computed in $\mathcal{V}$.
\end{remark}

\begin{proposition} \label{prop:riesz-map-isomorphism-criterion} \hlt{Riesz Map Isomorphism Criterion}\\
Let $\mathcal{V}$ be a metric vector space and define $\tau:\mathcal{V}\to\mathcal{V}'$ by $\tau(\mathbf{x})(\mathbf{v})=\innerprod{\mathbf{v},\mathbf{x}}$.
\begin{enumerate}[label=(\roman*)]
\setlength{\itemsep}{0pt}
\item $\tau\in\LL(\mathcal{V},\mathcal{V}')$.
\item $\ker(\tau)=\operatorname{rad}(\mathcal{V})$.
\item $\tau$ is an isomorphism $\Longleftrightarrow\operatorname{rad}(\mathcal{V})=\{\mathbf{0}\}$.
\end{enumerate}
\begin{proof}
Set $B(\mathbf{x},\mathbf{y})\defeq\innerprod{\mathbf{x},\mathbf{y}}$.
\begin{enumerate}[label=(\roman*)]
\setlength{\itemsep}{0pt}
\item[$(i)$] The identity $\tau=\Psi_B\in\LL(\mathcal{V},\mathcal{V}')$ holds by Proposition \ref{prop:bilinear-forms-linearize-through-duals}(i).
\item[$(ii)$] The identity $\ker(\tau)=\operatorname{rad}_{R}(B)$ holds by Proposition \ref{prop:bilinear-radicals-matrix-and-dual-maps}(i). Hence $\ker(\tau)=\operatorname{rad}(\mathcal{V})$ by Proposition \ref{prop:metric-radical-equals-bilinear-radicals}(i).
\item[$(iii)$] ($\Rightarrow$) If $\tau$ is an isomorphism, then it is injective by Proposition \ref{prop:linear-map-invertible-iff-bijective}.\\
Thus $\ker(\tau)=\{\mathbf{0}\}$ by Proposition \ref{prop:injective-iff-kernel-zero}. Hence $\operatorname{rad}(\mathcal{V})=\{\mathbf{0}\}$ by (ii).
\item[$(iii)$] ($\Leftarrow$) If $\operatorname{rad}(\mathcal{V})=\{\mathbf{0}\}$, then $\ker(\tau)=\{\mathbf{0}\}$ by (ii). Thus $\tau$ is injective by Proposition \ref{prop:injective-iff-kernel-zero}.\\
The equality $\dim(\mathcal{V})=\dim(\mathcal{V}')$ holds by Theorem \ref{thm:dual-basis-and-dimension}(ii).\\
Thus $\tau$ is surjective by Corollary \ref{cor:equal-dimension-kernel-range}(ii) and is an isomorphism by Proposition \ref{prop:linear-map-invertible-iff-bijective}.
\end{enumerate}
\hlt{Structural Commentary}: The radical is the Riesz-map kernel; equal dimensions convert injectivity into invertibility.\\
\hlt{Key Insight}: The radical is precisely the kernel of the Riesz map.
\end{proof}
\end{proposition}

\begin{theorem} \label{thm:riesz-map-for-nonsingular-metric-spaces} \hlt{Riesz Map for Nonsingular Metric Spaces}\\
Let $\mathcal{V}$ be a nonsingular metric vector space. Then $(\forall\,f\in\mathcal{V}')(\exists!\,\mathbf{x}\in\mathcal{V})(\forall\,\mathbf{v}\in\mathcal{V})\,f(\mathbf{v})=\innerprod{\mathbf{v},\mathbf{x}}$.
\begin{proof} Define $\tau:\mathcal{V}\to\mathcal{V}'$ by $\tau(\mathbf{x})(\mathbf{v})=\innerprod{\mathbf{v},\mathbf{x}}$. The map $\tau$ is an isomorphism by Proposition \ref{prop:riesz-map-isomorphism-criterion}(iii).\ Fix $f\in\mathcal{V}'$ and set $\mathbf{x}\defeq\tau^{-1}(f)$.\\
Then $(\forall\,\mathbf{v}\in\mathcal{V})\,f(\mathbf{v})=\tau(\mathbf{x})(\mathbf{v})=\innerprod{\mathbf{v},\mathbf{x}}$. If $\mathbf{y}\in\mathcal{V}$ satisfies $(\forall\,\mathbf{v}\in\mathcal{V})\,f(\mathbf{v})=\innerprod{\mathbf{v},\mathbf{y}}$, then $\tau(\mathbf{y})=f=\tau(\mathbf{x})$; $\mathbf{y}=\mathbf{x}$ by injectivity of $\tau$.\\
\hlt{Structural Commentary}: Riesz inversion constructs each representative; injectivity proves its uniqueness.\\
\hlt{Key Insight}: The radical is the sole obstruction to representing every functional.
\end{proof}
\end{theorem}

\begin{definition} \label{def:metric-isometry-and-classical-groups} \hlt{Isometries and Classical Groups}\\
Let $\mathcal{V},\mathcal{W}$ be metric vector spaces.
\begin{enumerate}[label=(\roman*)]
\setlength{\itemsep}{0pt}
\item An \hlt{isometry} $\sigma:\mathcal{V}\to\mathcal{W}$ is a bijective linear map satisfying $(\forall\,\mathbf{x},\mathbf{y}\in\mathcal{V})\,\innerprod{\sigma\mathbf{x},\sigma\mathbf{y}}=\innerprod{\mathbf{x},\mathbf{y}}$. Write $\mathcal{V}\cong\mathcal{W}$ if one exists.
\item The group of isometries $\mathcal{V}\to\mathcal{V}$ is $\operatorname{Iso}(\mathcal{V})$.\\
For nonsingular orthogonal $\mathcal{V}$, this is the \hlt{orthogonal group} $\mathcal{O}(\mathcal{V})$. For nonsingular symplectic $\mathcal{V}$, it is the \hlt{symplectic group} $\operatorname{Sp}(\mathcal{V})$.
\end{enumerate}
\end{definition}

\begin{proposition} \label{prop:matrix-criterion-for-metric-isometries} \hlt{Matrix Criterion for Isometries}\\
Let $\mathcal{V}$ be a metric vector space with $n\defeq\dim(\mathcal{V})\geq1$, and let $\beta$ be an ordered basis. Let $\sigma\in\LL(\mathcal{V})$, and set $M\defeq M_{\beta}$ and $A\defeq[\sigma]_{\beta\leftarrow\beta}$.
\begin{enumerate}[label=(\roman*)]
\setlength{\itemsep}{0pt}
\item $\sigma\in\operatorname{Iso}(\mathcal{V})\Longleftrightarrow\sigma$ is bijective and $A^tMA=M$.
\item If $\sigma\in\operatorname{Iso}(\mathcal{V})$ and $M$ is invertible, then $\det A=\pm1$.
\end{enumerate}
\begin{proof}
\begin{enumerate}[label=(\roman*)]
\setlength{\itemsep}{0pt}
\item[$(i)$] ($\Rightarrow$) Suppose $\sigma\in\operatorname{Iso}(\mathcal{V})$. Then $\sigma$ is bijective by definition.\\
The identity $(\forall\,X,Y\in\mathbb{F}^n)\,X^tA^tMAY=X^tMY$ holds by metric preservation and Proposition \ref{prop:metric-form-matrix-recovery-and-congruence}(i). Hence $A^tMA=M$ by evaluation at pairs of standard coordinate columns.
\item[$(i)$] ($\Leftarrow$) Suppose $\sigma$ is bijective and $A^tMA=M$.\\
Then $(\forall\,\mathbf{x},\mathbf{y}\in\mathcal{V})\,\innerprod{\sigma\mathbf{x},\sigma\mathbf{y}}=\innerprod{\mathbf{x},\mathbf{y}}$ by Proposition \ref{prop:metric-form-matrix-recovery-and-congruence}(i).\\
Hence $\sigma\in\operatorname{Iso}(\mathcal{V})$ by definition.
\item[$(ii)$] Suppose $\sigma\in\operatorname{Iso}(\mathcal{V})$ and $M$ is invertible. Then $A^tMA=M$ by (i).\\
Thus $\det(M)=\det(A^tMA)=(\det A)^2\det(M)$ by determinant multiplicativity and transpose invariance.\\
Thus $(\det A)^2=1$ by cancellation of $\det(M)\neq0$.\\
Hence $(\det A-1)(\det A+1)=0$, so $\det A\in\{1,-1\}$ by the no-zero-divisor property of $\mathbb{F}$.
\end{enumerate}
\hlt{Structural Commentary}: Coordinate evaluation identifies metric preservation with congruence; determinants constrain its stabilizer.\\
\hlt{Key Insight}: Isometries are exactly the invertible solutions of $A^tMA=M$.
\end{proof}
\end{proposition}

\begin{definition} \label{def:metric-adjoint-nonsingular-form} \hlt{Adjoint for a Nonsingular Bilinear Form}\\
Let $\mathcal{V}$ be nonsingular and $T\in\LL(\mathcal{V})$. An operator $S\in\LL(\mathcal{V})$ is a \hlt{metric adjoint} of $T$ $\Longleftrightarrow$ $(\forall\,\mathbf{x},\mathbf{y}\in\mathcal{V})\,\innerprod{T\mathbf{x},\mathbf{y}}=\innerprod{\mathbf{x},S\mathbf{y}}$.\\
If this operator is unique, denote it by $T^*$.
\end{definition}

\begin{proposition} \label{prop:metric-adjoints-and-isometries} \hlt{Metric Adjoints and Isometries}\\
Let $\mathcal{V}$ be nonsingular with form $B(\mathbf{x},\mathbf{y})=\innerprod{\mathbf{x},\mathbf{y}}$. Let $T\in\LL(\mathcal{V})$, and set $n\defeq\dim(\mathcal{V})$.
\begin{enumerate}[label=(\roman*)]
\setlength{\itemsep}{0pt}
\item $(\exists!\,T^*\in\LL(\mathcal{V}))(\forall\,\mathbf{x},\mathbf{y}\in\mathcal{V})\,\innerprod{T\mathbf{x},\mathbf{y}}=\innerprod{\mathbf{x},T^*\mathbf{y}}$.
\item If $n\geq1$ and $\beta$ is an ordered basis, set $M\defeq M_{\beta}$ and $A\defeq[T]_{\beta}$; then $[T^*]_{\beta}=M^{-1}A^tM$.
\item $T$ is an isometry $\Longleftrightarrow T$ is bijective and $T^*T=I$.
\item If $T$ is bijective, then $T^*T=I\Longleftrightarrow T^*=T^{-1}$.
\end{enumerate}
\begin{proof}
\begin{enumerate}[label=(\roman*)]
\setlength{\itemsep}{0pt}
\item[$(i)$] Define $\tau:\mathcal{V}\to\mathcal{V}'$ by $\tau(\mathbf{y})(\mathbf{x})=\innerprod{\mathbf{x},\mathbf{y}}$. Define $F:\mathcal{V}\to\mathcal{V}'$ by $F(\mathbf{y})(\mathbf{x})=\innerprod{T\mathbf{x},\mathbf{y}}$.\\
Then $\tau$ is an isomorphism by Proposition \ref{prop:riesz-map-isomorphism-criterion}(iii). The relation $F\in\LL(\mathcal{V},\mathcal{V}')$ holds by bilinearity.\\
Set $T^*\defeq\tau^{-1}\circ F\in\LL(\mathcal{V})$. Then $(\forall\,\mathbf{x},\mathbf{y}\in\mathcal{V})\,\innerprod{T\mathbf{x},\mathbf{y}}=\innerprod{\mathbf{x},T^*\mathbf{y}}$ by the definitions of $\tau$ and $F$.\\
If $S\in\LL(\mathcal{V})$ satisfies this identity, then $\tau\circ S=F=\tau\circ T^*$. Hence $S=T^*$ by injectivity of $\tau$.
\item[$(ii)$] Assume $n\geq1$. The matrix $M$ is invertible by Proposition \ref{prop:metric-radical-coordinate-criteria}(iii).\\
Set $D\defeq[T^*]_{\beta}$. Then $(\forall\,X,Y\in\mathbb{F}^n)\,X^tA^tMY=X^tMDY$ by the adjoint identity and Proposition \ref{prop:metric-form-matrix-recovery-and-congruence}(i).\\
Hence $A^tM=MD$ by evaluation at pairs of standard coordinate columns, and therefore $D=M^{-1}A^tM$.
\item[$(iii)$] ($\Rightarrow$) Suppose $T$ is an isometry. Then $T$ is bijective by definition.\\
For $\mathbf{x},\mathbf{y}\in\mathcal{V}$, $\innerprod{\mathbf{x},(T^*T-I)\mathbf{y}}=\innerprod{T\mathbf{x},T\mathbf{y}}-\innerprod{\mathbf{x},\mathbf{y}}$ by the adjoint identity.\\
This difference is $0$ by metric preservation, so $(T^*T-I)\mathbf{y}\in\operatorname{rad}_{R}(B)=\{\mathbf{0}\}$ by Proposition \ref{prop:metric-radical-equals-bilinear-radicals}(i); hence $T^*T=I$.
\item[$(iii)$] ($\Leftarrow$) Suppose $T$ is bijective and $T^*T=I$.\\
Then $(\forall\,\mathbf{x},\mathbf{y}\in\mathcal{V})\,\innerprod{T\mathbf{x},T\mathbf{y}}=\innerprod{\mathbf{x},T^*T\mathbf{y}}=\innerprod{\mathbf{x},\mathbf{y}}$ by the adjoint identity. Hence $T$ is an isometry by definition.
\item[$(iv)$] ($\Rightarrow$) If $T^*T=I$, then $T^*=T^{-1}$ by right-composition with $T^{-1}$.
\item[$(iv)$] ($\Leftarrow$) If $T^*=T^{-1}$, then $T^*T=T^{-1}T=I$.
\end{enumerate}
\hlt{Structural Commentary}: Riesz inversion constructs the adjoint; the radical test converts metric preservation into $T^*T=I$.\\
\hlt{Key Insight}: Nonsingularity turns the adjoint relation into a unique operator identity.
\end{proof}
\end{proposition}

\begin{proposition} \label{prop:metric-spaces-isometric-iff-congruent-matrices} \hlt{Isometry Is Matrix Congruence}\\
Let $(\mathcal{V},B_{\mathcal{V}})$ and $(\mathcal{W},B_{\mathcal{W}})$ be metric vector spaces with $\dim(\mathcal{V})=\dim(\mathcal{W})=n\geq1$.\\
Let $\beta,\gamma$ be ordered bases with respective form matrices $M_{\beta},N_{\gamma}$. Then $\mathcal{V}\cong\mathcal{W}\Longleftrightarrow M_{\beta}$ and $N_{\gamma}$ are congruent.
\begin{proof}
\begin{enumerate}[label=\roman*.]
\setlength{\itemsep}{0pt}
\item[$(\Rightarrow)$] Let $\sigma:\mathcal{V}\to\mathcal{W}$ be an isometry and set $A\defeq[\sigma]_{\gamma\leftarrow\beta}$. The map $\sigma$ is invertible by Proposition \ref{prop:linear-map-invertible-iff-bijective}.\\
By Proposition \ref{prop:composition-and-matrix-multiplication}, $[\sigma^{-1}]_{\beta\leftarrow\gamma}A=I_n=A[\sigma^{-1}]_{\beta\leftarrow\gamma}$, so $A$ is invertible.\\
The identity $(\forall\,X,Y\in\mathbb{F}^n)\,X^tM_{\beta}Y=X^tA^tN_{\gamma}AY$ holds by metric preservation and Proposition \ref{prop:metric-form-matrix-recovery-and-congruence}(i).\\
Hence $M_{\beta}=A^tN_{\gamma}A$ by evaluation at pairs of standard coordinate columns, so the matrices are congruent.
\item[$(\Leftarrow)$] Suppose $M_{\beta}$ and $N_{\gamma}$ are congruent. Fix an invertible $P\in\operatorname{Mat}(n,n,\mathbb{F})$ such that $N_{\gamma}=P^tM_{\beta}P$, and set $A\defeq P^{-1}$.\\
Then $M_{\beta}=A^tN_{\gamma}A$. Define $\sigma:\mathcal{V}\to\mathcal{W}$ by $[\sigma\mathbf{x}]_{\gamma}=A[\mathbf{x}]_{\beta}$. The map $\sigma$ is linear by Proposition \ref{prop:coordinate-map-isomorphism}(i) and Proposition \ref{prop:matrix-induced-map-linear}.\\
The coordinate rule with $A^{-1}$ defines its inverse, so $\sigma$ is bijective.\\
For $\mathbf{x},\mathbf{y}\in\mathcal{V}$ with coordinate columns $X,Y$, $B_{\mathcal{W}}(\sigma\mathbf{x},\sigma\mathbf{y})=X^tA^tN_{\gamma}AY$ by Proposition \ref{prop:metric-form-matrix-recovery-and-congruence}(i).\\
The identity $B_{\mathcal{V}}(\mathbf{x},\mathbf{y})=X^tM_{\beta}Y=X^tA^tN_{\gamma}AY$ holds by Proposition \ref{prop:metric-form-matrix-recovery-and-congruence}(i) and the congruence identity; $\sigma$ is an isometry.
\end{enumerate}
\hlt{Structural Commentary}: Coordinate matrices translate metric preservation into congruence and reconstruct the isometry.\\
\hlt{Key Insight}: Isometry classes are exactly congruence classes of form matrices.
\end{proof}
\end{proposition}

\begin{perspective} \label{per:ordinary-congruence-bilinear-form-bridge} \hlt{Congruence versus Similarity}\\
Similarity $N=P^{-1}MP$ governs operator matrices. Congruence $N=P^tMP$ governs bilinear-form matrices because the same basis change acts in both slots. Over $\mathbb{R}$, transpose equals conjugate transpose, so real form classification uses congruence.
\end{perspective}
